\section{How an audit is built}
\label{sec:architecture-pipeline}

\textbf{In brief.} An audit goes through six phases, each writing one document
that the next one reads. The application is the thinnest part: it accepts a
request, allows one run per instance, streams events and reads the artefacts back.
The engine is a pipeline that wraps every tool call in a governance boundary: a
call is checked against the budgets, archived with a reference, and archived even
when it is refused or fails. The policy only ever reasons over evidence that the
tools produced, and the score comes from a separate evaluator, not from the policy.

\textbf{What the code enforces, and what it does not.} What is enforced: a lost
piece of evidence stops the run; a deliverable that misses a device is rejected;
the report is assembled deterministically; the phase-five status can never be
certified. What is not guaranteed: the budgets only exist when they are passed
explicitly, the laboratory lock is optional, and two scenarios on a shared
laboratory are not mutually excluded.

\marginref{\code{docs/architecture/run-modes.md} for the modes; \code{src/agent/run\_modes.py} for the selection.}

\subsection{The application does not carry the logic}

The web layer validates the request, guarantees a single active run per instance,
streams the events and re-reads the artefacts. The strategy is chosen by
\code{make\_runner}, not by the interface: for \code{lance} and \code{rules} it
instantiates a pipeline and wraps it in an adapter that reuses the same object and
its public API; for \code{vanilla} and \code{scripted} it builds a standalone
runner that shares with the pipeline only the tool boundary, the evidence log, the
cost tracking and the laboratory reservation.

\begin{table}[H]
\centering
\caption{Execution modes and their boundary.}
\label{tab:run-modes}
\small
\begin{tabular}{@{}llll@{}}
\toprule
Mode & Engine & Model & Effects \\
\midrule
\code{lance} & full pipeline, six phases & required & scenario, deployment, evaluation \\
\code{rules} & pipeline, phases 3--4 & none & scenario and deterministic audit \\
\code{vanilla} & standalone runner & required & bounded network probe on an address range \\
\code{scripted} & standalone runner & none & versioned procedures \\
\bottomrule
\end{tabular}
\end{table}

The \code{rules} mode is not a second engine: the factory sets the decision policy,
restricts the phases to analysis and verification, and refuses any provider, any
direct network access and any monetary budget. The two standalone modes, in turn,
refuse a scenario and every harness-specific setting. The qualifier
\code{full}/\code{compact} designates the orchestration, never the type of runner.

\subsection{The six phases and what each hands over}

The registry assigns to each phase a prompt, a tool group, a validator and a
deliverable; the routing is an explicit table, not a dynamic import.

\begin{table}[H]
\centering
\caption{Phases, main deliverables and associated artefacts.}
\label{tab:phases}
\small
\begin{tabular}{@{}cllp{4.4cm}@{}}
\toprule
Phase & Main deliverable & Who decides & Associated artefacts \\
\midrule
1 & \code{01\_graph\_analysis.md} & code (projection of the graph tools) & \code{01\_graph\_evidence.json} \\
2 & \code{02\_recon.md} & the model chooses the probes & \code{02\_recon\_evidence.json}, \code{02\_topology\_edges.json} \\
3 & \code{03\_vuln\_analysis.json} & the model writes the hypotheses & \code{03\_vuln\_analysis\_raw.json}, \code{03\_phase3\_status.json}, per-device sheets \\
4 & \code{04\_exploitation.json} & the model tests one flaw per micro-agent & \code{execution\_errors} diagnostics \\
5 & \code{05\_intrusion.json} & code (assessment rendered from the log) & \code{05\_intrusion\_context.json} \\
6 & \code{06\_report.md} & code (assembled from section sheets) & \code{06\_report\_analysis.md} \\
\bottomrule
\end{tabular}
\end{table}

Phase 1 projects the results of the graph tools into a log of facts independent of
the model's discourse, and the runner rejects a submission whose per-device
coverage is incomplete. Phase 2 lets the model choose the probes, but the recon
document is rendered from the tool log rather than copied by the model. Phase 3
runs in three steps: deterministic scanning, one analyst per device, then an
aggregation that normalises, deduplicates and projects the findings. Phase 4
launches one micro-agent per vulnerability. Phase 5 produces its assessment from
the log: the model supplies an end marker only, and the status remains
\code{not\_certified}. Phase 6 is the extreme case: its agent declares no tool and
runs a single turn, so that the report is assembled and the model writes only a
bounded analyst note.

\marginref{Phase modules: \code{src/agent/phases/registry.py}, \code{PHASE\_MODULES}. Phase-1 projection: \code{GraphPhase.\_build\_graph\_evidence\_projection}.}

\subsection{Graph and run context}

The graph is built in two layers: a loader reads the infrastructure description
into data objects, then populates a backend. The backend is an abstract interface
with a network-graph implementation, so that another engine can be substituted
without touching the tools. Two details matter for reading results: the neighbour
query is the sorted union of successors and predecessors, and path search uses an
undirected view so that a path is found in both directions.

The state belongs to the run, not to the process. Each pipeline creates an
immutable run context carrying the backend, the infrastructure, the vulnerability
reports, the risk scores, the scenario or discovery topology and the weighted
graph. The graph tools retrieve that context without changing the arguments
exposed to the model; thread workers capture the owner explicitly rather than
relying on implicit propagation.

\marginref{\code{src/loader.py} and \code{src/graph\_backend.py}; context in \code{src/agent/core/graph\_context.py}.}

\subsection{The tool governance boundary}

All profiles and all modes go through the same wrapper. On each call it checks the
previous piece of evidence, the stop request and the budgets; increments a counter
and refuses the call if the ceiling is reached; assigns a reference; applies the
scope refusals; executes; and archives an entry, including for a refusal or an
exception.

\begin{table}[H]
\centering
\caption{Controls applied to every tool call and the consequence of a failure.}
\label{tab:executor}
\small
\begin{tabular}{@{}p{3.6cm}p{5.2cm}p{5.0cm}@{}}
\toprule
Control & Mechanism & Consequence of a failure \\
\midrule
Evidence integrity & a flag set when the archive cannot be written & the run stops, and is not verifiable \\
User stop & a stop event, in the experimental scope & the run stops \\
Time budget & monotonic clock against the limit & the run stops \\
Monetary budget & the cost tracker, on observed usage & the run stops \\
Call budget & a counter taken under the archive lock & the run stops \\
Scope & the destination fields against the run's target network & the call is refused and archived \\
Free execution & the execution tool & refused as unverifiable scope \\
\bottomrule
\end{tabular}
\end{table}

Two properties matter for interpreting a result. A failure to write the log is
never ignored: it interrupts the run rather than leaving untraced actions, which
is the difference between an audit and a log that only records what went well. And
the cost limit applies to observed usage, not to a reservation of calls still in
flight, so a run can exceed its monetary budget by the cost of the calls it had
already started.

\marginref{Implementation: \code{wrap\_tool} and \code{check\_execution\_limits}, \code{src/agent/core/executor.py}.}

\subsection{Artefacts, cost tracking and retention}

The run folder is created exclusively under the output tree; in the event of a
collision a hexadecimal suffix is appended and no earlier folder is overwritten.
The run metadata records the commit, the policies and the limits, and the cost
summary is written in the finalisation block, hence even when the run raises. A
submission bundle is a zip equipped with a SHA-256 manifest that is re-checked
after writing; the download route excludes private files and hidden entries.

The price of a phase is resolved once and recorded with its source, so that a
completed run cannot be revalued after the fact; the summary distinguishes the
cumulative duration of the agents from the wall-clock time, which are not the same
measure.

\marginref{Collection contract: \code{src/benchmark/artifacts.py}. Pricing: \code{src/agent/cost\_tracker.py}.}

\subsection{What the architecture does not guarantee}

The isolation is per run for the deliverables and the graph, but not entirely so:
the phase services may still call the common engine, and the simultaneity of two
scenarios on a shared laboratory is not guaranteed. The laboratory reservation is
active only when its variable is set, and is reentrant only within one thread.
Without explicit parameters, the interactive pipeline imposes neither a call
ceiling nor an analysis timeout per device, so the budgets of
table~\ref{tab:executor} are conditional rather than structural.

Three points could not be established by reading the code, and are therefore not
asserted: no archiving or compaction mechanism for run folders was found, so
retention appears to be that of the file system; one evidence file is read by two
modules but its main producer was not identified; and the exact field content of
the provider diagnostics log was not verified line by line. Finally, a completed
status means that the runner finished its work, never that the target was
exhaustively covered nor that an access was proven.
